\documentclass[10pt,a4paper]{amsart}
\usepackage{amsmath,amssymb,amsthm,graphicx}
% Retained minimal preamble for the Stacks Project chapter flat.tex (More on Flatness).
% Verbatim excerpt of the chapter preamble: Xy-pic setup (native preamble lines 14-18) and the native
% theorem declarations with the mathematical macros (native preamble lines 154-217). Copied mechanically;
% nothing added, removed or redefined. The wrapper class and the amsmath/amssymb/amsthm/graphicx base
% packages are supplied by the assembler, so the class/externaldocument part of the native preamble is omitted.
% For commutative diagrams we use Xy-pic
\usepackage[all]{xy}

% We use 2cell for 2-commutative diagrams.
\xyoption{2cell}
% Theorem environments.
%
\theoremstyle{plain}
\newtheorem{theorem}[subsection]{Theorem}
\newtheorem{proposition}[subsection]{Proposition}
\newtheorem{lemma}[subsection]{Lemma}

\theoremstyle{definition}
\newtheorem{definition}[subsection]{Definition}
\newtheorem{example}[subsection]{Example}
\newtheorem{exercise}[subsection]{Exercise}
\newtheorem{situation}[subsection]{Situation}

\theoremstyle{remark}
\newtheorem{remark}[subsection]{Remark}
\newtheorem{remarks}[subsection]{Remarks}

\numberwithin{equation}{subsection}

% Macros
%
\def\lim{\mathop{\mathrm{lim}}\nolimits}
\def\colim{\mathop{\mathrm{colim}}\nolimits}
\def\Spec{\mathop{\mathrm{Spec}}}
\def\Hom{\mathop{\mathrm{Hom}}\nolimits}
\def\Ext{\mathop{\mathrm{Ext}}\nolimits}
\def\SheafHom{\mathop{\mathcal{H}\!\mathit{om}}\nolimits}
\def\SheafExt{\mathop{\mathcal{E}\!\mathit{xt}}\nolimits}
\def\Sch{\mathit{Sch}}
\def\Mor{\mathop{\mathrm{Mor}}\nolimits}
\def\Ob{\mathop{\mathrm{Ob}}\nolimits}
\def\Sh{\mathop{\mathit{Sh}}\nolimits}
\def\NL{\mathop{N\!L}\nolimits}
\def\CH{\mathop{\mathrm{CH}}\nolimits}
\def\proetale{{pro\text{-}\acute{e}tale}}
\def\etale{{\acute{e}tale}}
\def\QCoh{\mathit{QCoh}}
\def\Ker{\mathop{\mathrm{Ker}}}
\def\Im{\mathop{\mathrm{Im}}}
\def\Coker{\mathop{\mathrm{Coker}}}
\def\Coim{\mathop{\mathrm{Coim}}}

% Boxtimes
%
\DeclareMathSymbol{\boxtimes}{\mathbin}{AMSa}{"02}

%
% Macros for moduli stacks/spaces
%
\def\QCohstack{\mathcal{QC}\!\mathit{oh}}
\def\Cohstack{\mathcal{C}\!\mathit{oh}}
\def\Spacesstack{\mathcal{S}\!\mathit{paces}}
\def\Quotfunctor{\mathrm{Quot}}
\def\Hilbfunctor{\mathrm{Hilb}}
\def\Curvesstack{\mathcal{C}\!\mathit{urves}}
\def\Polarizedstack{\mathcal{P}\!\mathit{olarized}}
\def\Complexesstack{\mathcal{C}\!\mathit{omplexes}}
% \Pic is the operator that assigns to X its picard group, usage \Pic(X)
% \Picardstack_{X/B} denotes the Picard stack of X over B
% \Picardfunctor_{X/B} denotes the Picard functor of X over B
\def\Pic{\mathop{\mathrm{Pic}}\nolimits}
\def\Picardstack{\mathcal{P}\!\mathit{ic}}
\def\Picardfunctor{\mathrm{Pic}}
\def\Deformationcategory{\mathcal{D}\!\mathit{ef}}

% --- appended 2026-10-12 (Chief): the project own definitions of the reference/slogan/history
% display environments, copied mechanically from the shared Stacks preamble lines 8-12
% (\usepackage{verbatim} plus the three \newenvironment lines, which the chapter preamble keeps
% verbatim from the project preamble). Nothing added, removed or redefined: these environments
% carry no mathematical content and the project itself typesets them as comments. Without these
% three lines a counted statement whose own text uses them fails with Environment slogan undefined.
\usepackage{verbatim}
\newenvironment{reference}{\comment}{\endcomment}
\newenvironment{slogan}{\comment}{\endcomment}
\newenvironment{history}{\comment}{\endcomment}
\usepackage{cleveref}
\crefname{originallemma}{original lemma}{original lemmas}
\Crefname{originallemma}{Original lemma}{Original lemmas}
\makeatletter
\newlabel{lemma-fitting-ideals-complex}{{0ESP}{source}{}{sourcepointer-1}{}}
\newlabel{lemma-fitting-ideals-complex@cref}{{[originallemma][0][]0ESP}{[1][0][]source}{}{}{}}
\makeatother
\crefname{originallemma}{original lemma}{original lemmas}
\Crefname{originallemma}{Original lemma}{Original lemmas}
\makeatletter
\newlabel{divisors-lemma-blowing-up-two-ideals}{{080A}{source}{}{sourcepointer-2}{}}
\newlabel{divisors-lemma-blowing-up-two-ideals@cref}{{[originallemma][0][]080A}{[1][0][]source}{}{}{}}
\makeatother
\crefname{originallemma}{original lemma}{original lemmas}
\Crefname{originallemma}{Original lemma}{Original lemmas}
\makeatletter
\newlabel{divisors-lemma-blowup-fitting-ideal}{{0ESN}{source}{}{sourcepointer-3}{}}
\newlabel{divisors-lemma-blowup-fitting-ideal@cref}{{[originallemma][0][]0ESN}{[1][0][]source}{}{}{}}
\makeatother
\crefname{originallemma}{original lemma}{original lemmas}
\Crefname{originallemma}{Original lemma}{Original lemmas}
\makeatletter
\newlabel{divisors-lemma-dominate-admissible-blowups}{{080N}{source}{}{sourcepointer-4}{}}
\newlabel{divisors-lemma-dominate-admissible-blowups@cref}{{[originallemma][0][]080N}{[1][0][]source}{}{}{}}
\makeatother
\title{A perfect complex whose cohomology is locally free on a dense open becomes of tor dimension at most one after a suitable admissible blowup}
\author{The Stacks Project Authors}
\date{stacks-project:flat.tex:local-snapshot-20261011; GNU FDL 1.2 or later}
\begin{document}
\maketitle
\noindent\textit{Editorial note.} All delivered mathematical text is the author's own, copied from the pinned Stacks Project chapter file flat.tex; no statement or proof step was rewritten, reordered or removed.\par\medskip
\noindent\textit{Editorial note.} The preamble is the chapter's own retained minimal preamble (sources/stacks-flat/retained-preamble.tex), a verbatim excerpt of the native Xy-pic setup and of the native theorem declarations and mathematical macros, followed by the project own definitions of the reference, slogan and history display environments (verbatim package plus three newenvironment lines) that the retained file carries at its end and that the project itself typesets as comments. No macro was invented or redefined, and the wrapper class and the amsmath/amssymb/amsthm/graphicx base packages are supplied by the assembler.\par\medskip
\noindent\textit{Editorial note.} The counted statement and proof contain no cross reference, no citation and no included file, so the delivered body is self contained.\par\medskip
\noindent\textit{Editorial note.} Canonical location of the counted proposition: Stacks tag 0ESQ (sources/stacks-flat/pinned-original-source.tex line 13036).\par\medskip
\noindent\textit{Allowed standard prerequisites (editorial).}\begin{itemize}
\item The global Fitting ideals Fit\_\{i,k\}(E) of a perfect complex and their independence of the representative. Same chapter, Lemma labelled lemma-fitting-ideals-complex, cited at proof lines 13067-13070. Identifies the ideal of minors with a globally defined ideal of O\_X.
\item Blowing up in a product of ideals and admissibility of the blowup in an invertible Fitting ideal. Divisors, Lemma labelled divisors-lemma-blowing-up-two-ideals, cited at proof lines 13079-13081, and Lemma labelled divisors-lemma-blowup-fitting-ideal, cited at proof lines 13081-13083 and 13092-13094. Shows that the product blowup dominates the individual admissible blowups and controls the tor dimension after the blowup.
\item Domination of admissible blowups. Divisors, Lemma labelled divisors-lemma-dominate-admissible-blowups, cited at proof lines 13083-13085. Supplies the comparison with the proof that the dominating blowup is again U-admissible.
\end{itemize}\medskip
\ifdefined\hypertarget\hypertarget{sourcepointer-1}{}\fi
\noindent\textit{Source reference.} Original lemma 0ESP: sources/stacks-flat/pinned-original-source.tex, source lines 12963--12963. Same-chapter pointer into the author environment lemma at source line 12963; invoked at proof line 13066 of the counted proof. It is retained only as a navigation pointer; every load-bearing step is inside the delivered span.\par\medskip
\ifdefined\hypertarget\hypertarget{sourcepointer-2}{}\fi
\noindent\textit{Source reference.} Original lemma 080A: sources/stacks-chow/tags.txt, source lines 10175--10175. Cross-chapter standard prerequisite located in the project tag list staged in the local snapshot; invoked at proof line 13079 of the counted proof. It is retained only as a navigation pointer; every load-bearing step is inside the delivered span.\par\medskip
\ifdefined\hypertarget\hypertarget{sourcepointer-3}{}\fi
\noindent\textit{Source reference.} Original lemma 0ESN: sources/stacks-chow/tags.txt, source lines 18488--18488. Cross-chapter standard prerequisite located in the project tag list staged in the local snapshot; invoked at proof line 13080 of the counted proof. It is retained only as a navigation pointer; every load-bearing step is inside the delivered span.\par\medskip
\ifdefined\hypertarget\hypertarget{sourcepointer-4}{}\fi
\noindent\textit{Source reference.} Original lemma 080N: sources/stacks-chow/tags.txt, source lines 10188--10188. Cross-chapter standard prerequisite located in the project tag list staged in the local snapshot; invoked at proof line 13083 of the counted proof. It is retained only as a navigation pointer; every load-bearing step is inside the delivered span.\par\medskip
\begin{lemma}
\label{lemma-blowup-complex}
Let $X$ be a scheme. Let $E \in D(\mathcal{O}_X)$ be perfect.
Let $U \subset X$ be a scheme theoretically dense open subscheme
such that $H^i(E|_U)$ is finite locally free of constant rank $r_i$
for all $i \in \mathbf{Z}$.
Then there exists a $U$-admissible blowup $b : X' \to X$ such that
$H^i(Lb^*E)$ is a perfect $\mathcal{O}_{X'}$-module
of tor dimension $\leq 1$ for all $i \in \mathbf{Z}$.
\end{lemma}

\begin{proof}
We will construct and study the blowup affine locally. Namely, suppose that
$V \subset X$ is an affine open subscheme such that
$E|_V$ can be represented by the complex
$$
\mathcal{O}_V^{\oplus n_a} \xrightarrow{d_a}
\ldots \xrightarrow{d_{b - 1}} \mathcal{O}_V^{\oplus n_b}
$$
Set $k_ i = r_{i + 1} - r_{i + 2} + \ldots + (-1)^{b - i + 1}r_b$.
A computation which we omit show that over $U \cap V$ the rank of
$d_i$ is
$$
\rho_i = - k_i + n_{i + 1} - n_{i + 2} + \ldots + (-1)^{b - i + 1}n_b
$$
in the sense that the cokernel of $d_i$ is finite locally
free of rank $n_{i + 1} - \rho_i$. Let
$\mathcal{I}_i \subset \mathcal{O}_V$ be the ideal generated by the minors
of size $\rho_i \times \rho_i$ in the matrix of $d_i$.

\medskip\noindent
On the one hand, comparing with Lemma \ref{lemma-fitting-ideals-complex}
we see the ideal $\mathcal{I}_i$ corresponds to the global ideal
$\text{Fit}_{i, k_i}(E)$ which was shown to be
independent of the choice of the complex representing $E|_V$.
On the other hand, $\mathcal{I}_i$ is the $(n_{i + 1} - \rho_i)$th
Fitting ideal of $\Coker(d_i)$. Please keep this in mind.

\medskip\noindent
We let $b : X' \to X$ be the blowing up in the product of the ideals
$\text{Fit}_{i, k_i}(E)$; this makes sense as locally on $X$
almost all of these ideals are equal to the unit ideal (see above).
This blowup dominates the blowups $b_i : X'_i \to X$ in
the ideals $\text{Fit}_{i, k_i}(E)$, see
Divisors, Lemma \ref{divisors-lemma-blowing-up-two-ideals}.
By Divisors, Lemma \ref{divisors-lemma-blowup-fitting-ideal}
each $b_i$ is a $U$-admissible blowup. It follows that
$b$ is a $U$-admissible blowup (tiny detail omitted; compare with
the proof of Divisors, Lemma \ref{divisors-lemma-dominate-admissible-blowups}).
Finally, $U$ is still a scheme theoretically dense open subscheme of $X'$.
Thus after replacing $X$ by $X'$ we end up in the situation
discussed in the next paragraph.

\medskip\noindent
Assume $\text{Fit}_{i, k_i}(E)$ is an invertible ideal for all $i$.
Choose an affine open $V$ and a complex of finite free modules
representing $E|_V$ as above. It follows from
Divisors, Lemma \ref{divisors-lemma-blowup-fitting-ideal}
that $\Coker(d_i)$ has tor dimension $\leq 1$. Thus
$\Im(d_i)$ is finite locally free as the kernel of a map
from a finite locally free module to a finitely presented
module of tor dimension $\leq 1$. Hence $\Ker(d_i)$
is finite locally free as well (same argument). Thus the short
exact sequence
$$
0 \to \Im(d_{i - 1}) \to \Ker(d_i) \to H^i(E)|_V \to 0
$$
shows what we want and the proof is complete.
\end{proof}

\end{document}
